\documentclass[11pt]{article}
\usepackage[margin=1in]{geometry}
\usepackage{amsmath,amssymb,amsthm,booktabs,hyperref}
\setlength{\emergencystretch}{2em}
\newtheorem{theorem}{Theorem}
\newtheorem{lemma}{Lemma}
\title{A rank-one counterexample to conjecture 00000001101}
\author{ziangni-sys}
\date{4 October 2026}
\begin{document}
\maketitle
\begin{abstract}
The root system $A_1$ disproves the claim that every root system has a
Kazhdan--Lusztig polynomial whose leading coefficient has an odd prime factor.
Its three nonzero Kazhdan--Lusztig polynomials are all the constant polynomial
$1$. A self-contained Lean 4 project proves the conclusion for every
polynomial family satisfying the standard rank-one KL conditions, rather
than merely checking a supplied table.
\end{abstract}

\section{The statement and counterexample}
The original English statement is:
\begin{quote}
Definition: The Kazhdan--Lusztig polynomials $P_{x,w}(q)$ of Weyl groups.
Conjecture: For every root system there exist $x,w$ such that the leading
coefficient of $P_{x,w}$ contains an odd prime factor beyond powers of two
(diversity of prime factors of KL coefficients).
\end{quote}
The universal quantifier includes the reduced, irreducible, crystallographic
root system $\Phi=\{-1,1\}\subset\mathbb R$. The Euclidean inner product
is $\langle a,b\rangle=ab$. For either root $\alpha$, the root reflection
is
\[
r_\alpha(t)=t-2\frac{\langle t,\alpha\rangle}{\langle\alpha,\alpha\rangle}
\alpha=-t.
\]
Thus $\Phi$ spans its one-dimensional ambient space, is stable under
root reflections, satisfies the reducedness condition, and its Cartan
integers $2\langle\beta,\alpha\rangle/\langle\alpha,\alpha\rangle$
are $\pm2$. This is the standard root system $A_1$. Its Weyl group is
$W=\{e,s\}$ with $s^2=e$. The Coxeter lengths are $\ell(e)=0$ and
$\ell(s)=1$; its Bruhat order is $e<s$.

\section{Derivation of the KL polynomials}
We use the usual convention $P_{x,w}=0$ when $x\not\leq w$. For comparable
pairs the KL polynomials obey
\[
P_{w,w}=1,\qquad P_{x,w}(0)=1,\qquad
\deg P_{x,w}\leq\frac{\ell(w)-\ell(x)-1}{2}\quad(x<w).
\]
For $e<s$, the degree bound is zero, so $P_{e,s}$ is constant;
its constant term forces $P_{e,s}=1$. Diagonal normalization gives the
other two nonzero polynomials. Every pair in $W\times W$ is accounted for:
\[
\begin{array}{c|cc}
P_{x,w}&w=e&w=s\\\hline
x=e&1&1\\
x=s&0&1
\end{array}
\]

For completeness, the rank-one calculation can also be obtained directly
from the defining Hecke-algebra construction, without using a precomputed
table. Set $q=v^2$ and let $T=T_s$. The relation is
\[
T^2=(q-1)T+q,
\quad T^{-1}=q^{-1}T+(q^{-1}-1).
\]
The bar operation sends $v$ to $v^{-1}$ and $T$ to $T^{-1}$. Consequently
\[
\overline{v^{-1}(T+1)}
=v\bigl(q^{-1}T+(q^{-1}-1)+1\bigr)
=v^{-1}(T+1).
\]
This element has the prescribed leading term $v^{-1}T$ and lower
coefficient $v^{-1}\in v^{-1}\mathbb Z[v^{-1}]$. It is therefore the
canonical basis element $C'_s$. Indeed, the difference of two such
bar-invariant elements would be a scalar Laurent polynomial supported
in strictly negative powers of $v$ and equal to its bar, hence zero.
Comparing
\[
C'_s=v^{-\ell(s)}\sum_{x\leq s}P_{x,s}(v^2)T_x
=v^{-1}(T+1)
\]
again yields $P_{e,s}=P_{s,s}=1$.

\begin{theorem}
Conjecture 00000001101 is false.
\end{theorem}
\begin{proof}
Every nonzero KL polynomial of the root system $A_1$ equals $1$, so its
leading coefficient is $1$. No prime divides $1$: if a natural prime $p$
divided $1$, then $p=1$, contradicting $p\geq2$. In particular no odd
prime divides the leading coefficient. The only remaining pair $(s,e)$
has the zero polynomial, which has no leading nonzero coefficient and
no prime factorization. It cannot witness the conjecture. The existential
claim fails for this root system, disproving the universal statement.
\end{proof}

\section{Formal verification and semantic correspondence}
The project uses Lean 4.19.0 and its standard library only.
\texttt{W} is the two-element Weyl group indexing type, \texttt{action}
is identity or integer negation, and \texttt{reflection\_preserves\_roots}
checks preservation of the coordinate representatives $\{-1,1\}$.
\texttt{len} and \texttt{bruhat} encode the lengths and order above.
\texttt{Poly} encodes an integer polynomial as a coefficient function
$\mathbb N\to\mathbb Z$ together with a finite support bound.

\texttt{KLFamily P} is an explicit proposition on an arbitrary family
of these polynomials. It records diagonal normalization, zero outside
Bruhat order, constant term $1$ for comparable pairs, and the strict
degree bound coefficientwise: for $x<w$, a coefficient at $k$ vanishes
if $\ell(w)-\ell(x)\leq2k$. This is exactly
$2\deg P_{x,w}<\ell(w)-\ell(x)$ for nonzero polynomials.
In rank one these conditions uniquely determine all coefficients;
\texttt{KL\_unique\_coeff} proves that fact. Thus the theorem applies
to the actual KL family, which satisfies these standard conditions,
and to any proposed implementation of it.

\texttt{IsLeading p k} requires the coefficient at $k$ to be nonzero
and all higher coefficients to vanish. This faithfully excludes the
zero polynomial. \texttt{Prime} uses $p\geq2$ and the usual positive
natural divisor condition; \texttt{OddPrimeDivisor} requires primality,
oddness and divisibility of the absolute integer coefficient.
\texttt{leading\_is\_one} proves the leading coefficient is $1$ for
every pair and every proposed leading index.
\texttt{rank\_one\_disproof} then proves the negation of the existential
claim for every \texttt{KLFamily}; \texttt{certified\_counterexample}
instantiates it with a family whose defining conditions are themselves
proved in \texttt{table\_is\_KL}.

No KL property is declared as a Lean axiom. The properties are explicit
hypotheses of the generic theorem and are discharged for the rank-one
witness. The project does not formalize general root-system or Hecke-algebra
theory; the identification with $A_1$ and the direct canonical-basis
derivation are provided above. No general KL positivity theorem or
numerical search is needed. The final theorems depend only on the standard
logical axioms \texttt{propext} and \texttt{Quot.sound}; there are no
admitted proofs, custom axioms or native evaluation tactics.

\section{Reproduction and sources}
Run \texttt{lake build} and \texttt{lake env lean Main.lean} inside
\texttt{lean/}. Run \texttt{python verify.py} in the submission directory
for an independent exact Laurent-polynomial verification of the bar
calculation and all four table entries. Compile this source with
\texttt{tectonic report.tex}. The accompanying verification log records
the compiler results and visual PDF inspection.

The mathematical conventions are those of D.~Kazhdan and G.~Lusztig,
\emph{Representations of Coxeter groups and Hecke algebras},
Inventiones Mathematicae 53 (1979), 165--184,
\href{https://doi.org/10.1007/BF01390031}{doi:10.1007/BF01390031}.
For coefficient and indexing conventions see M.~Goresky's
\href{https://www.math.ias.edu/~goresky/tables.html}{Tables of
Kazhdan--Lusztig Polynomials} (Institute for Advanced Study).
The repository statement is
\href{https://github.com/The-Last-Math-Competition/The-Last-Math-Competition/blob/main/conjectures/00000001101.md}{conjectures/00000001101.md}.
\end{document}
